\odiseachapter{arco}{El arco funesto}\label{cap:arco}


La palabra es \gr{μνηστηροφονία} (mnēstērophonía), derivada de \gr{μνηστήρ} (pretendiente) y \gr{φονία} (matanza, homicidio). El griego solo necesita un vocablo donde el castellano precisa de cinco.

Un solo vocablo para nombrar el asesinato a sangre fría de ciento ocho jóvenes, seguido de una implacable ejecución.

Nolan y todos los directores de cine anteriores a él se las han tenido que ver con este dato. Cierto, los pretendientes eran unos gorrones, y no es menos cierto que algunos de ellos eran enemigos declarados del héroe. Antínoo es un miserable que intenta emboscar a Telémaco, involucrando a otros veinte secuaces. Eurímaco y Ctesipo pagan cara su arrogancia. Agelao intenta organizar a los aterrorizados galanes y se comprende que Odiseo y su gente los eliminen. Pero todos sumados no llegan a treinta, lo que quiere decir que los otros ochenta hombres son masacrados como ganado, sin necesidad y sin piedad alguna.

No hay forma de que Hollywood muestre semejante salvajada. Pero tampoco era necesario endulzar la historia hasta hacerla irreconocible, empezando por el caso de Telémaco. En su \emph{Nodisea}, el versionante director insiste en tratar al hijo de Ulises, que se parece a su padre en muchos aspectos ---incluyendo la crueldad---, como un muchachito atribulado e indeciso, además de físicamente torpe.

El resultado es ridículo y arruina una película que cuenta, como hemos visto, con hallazgos excelentes. Cuando los pretendientes intentan emboscar a Telémaco en el templo de Pilos (una escena inventada que no aparece en la épica de Homero, donde el príncipe, avisado por la metomentodo Atenea, simplemente da esquinazo a sus enemigos), tiene que ser papá quien intervenga para proteger al chaval. Peor aún, durante la matanza, \emph{daddy} le ordena que se esté quietecito, no se vaya a manchar las manos de sangre ---da risa imaginarse a Odiseo, rodeado de cien enemigos, preocupado por el bienestar emocional de su mocito--- y a Chris Nodoyuna no le queda otra que mandarlo al piso de arriba, a darse bofetadas con el traidor cabrero Melantio, con el único ánimo de tenerlo ocupado.

No menos caricaturesca es la inverosímil batalla en la que Matt Damon acaba hecho un San Sebastián, a pesar de lo cual Nolan no consigue capturar ni la violencia de la masacre ni el salvajismo del héroe. Es más, cuando parte de los pretendientes «hinca la rodilla» y se rinde, Ulises les perdona generosamente y con ello nos dan gato por liebre, o, más precisamente, gato por tigre. El Odiseo que nos muestra el canto XXII es exactamente eso: un despiadado tigre de Bengala, decidido a devorar hasta el último de sus enemigos. Convertirlo en un gato de Angora es, se mire como se mire, imperdonable.


Pero dejemos que sea Homero quien nos cuente la historia antes de seguir con el delicado asunto de cómo Hollywood adultera el clásico. En el canto XXI, Penélope invita a los pretendientes a la prueba del célebre arco:

\begin{versocita}
«Oídme, altivos galanos, que habéis tomado este techo\\
y un día y otro sin fin estáis comiendo y bebiendo\\
de un hombre que hace ya mucho que falta de él; y pretexto\\
otro no habéis encontrado a decir de palabra que esto:\\
que ansia tenéis por llevarme de esposa y en casamiento.\\
Pues, venga, mis pretendientes, aquí a la vista está el premio.\\
Os dejaré este gran arco del biendivino Odiseo;\\
con el que tense en sus manos el arco en menos esfuerzo\\
y por las hachas, las doce, la flecha la pase en un vuelo,\\
con él me iré aquí dejando mi casa de casamiento,\\
este casal tan hermoso y tan de vivires lleno\\
del que me voy a acordar me creo que hasta en los sueños».
\end{versocita}

Toda la escena está muy bien conseguida en el film de Nolan, aunque se salta un detalle esencial. En la \emph{Odisea}, Telémaco es el primero en intentar tensar el arco:

\begin{versocita}
[...] y soltó de sus hombros el manto purpurigrana\\
brincando en pie, y de sus hombros soltó la espada aguzada.\\
Hincó primero las hachas abriéndoles antes por cama\\
un largo surco, y dejábalas todas a cuerda alineadas\\
apelmazando la tierra en redor; con asombro miraban\\
lo bien a ley que ordenó ---y antes nunca él lo viera--- las hachas.\\
Y fue a probarse en el arco tras apostarse a la entrada.\\
Tres veces lo hizo temblar, de llevarlo a la gaza en gran ansia,\\
tres de su fuerza falló, esperando las tres en su alma\\
tensar el nervio y pasar por el hierro la flecha lanzada.\\
Y bien lo habría encordado tirando en más fuerza a la cuarta,\\
mas Odiseo de seña negó y reprimiose él la gana.
\end{versocita}

Observa, asombrada lectora, perplejo lector, lo que nos dice el Poeta. \emph{Telémaco es capaz de tensar el arco de Odiseo} ---al cuarto intento---, es decir, \emph{es casi tan fuerte como su padre} y, de hecho, casi tan astuto como este. Reprime sus ganas de hacerlo y dice así:

\begin{versocita}
«Áyayme, ya un parapoco, un indigno seré, bien lo creo,\\
o a lo mejor soy muy joven y manos aún no las tengo\\
para guardarme de quien me llegó en malhechuras primero.\\
Venid vosotros, que en fuerzas muy por delante yo os veo,\\
haced del arco el intento, y que se cumpla ya el juego».
\end{versocita}

El primero de los jóvenes que prueba es Leodes, que, entre los pretendientes, oficia como \gr{θυοσκόος} (thyoskóos), un término compuesto a partir de \gr{θύος}, «ofrenda quemada, incienso», y una raíz que significa «observar». Es decir, Leodes es «el que observa las ofrendas» o el adivino del grupo. Además es el más decente de la banda de gorrones:

\begin{versocita}
Leodes, hijo de Énope en pie se puso el primero;\\
él les hacía de arúspice y junto a la crátera asiento\\
tomaba siempre al rincón; y era el único allí desafecto\\
a sus excesos, y andaba indignado con todo el cortejo;\\
así que el arco y la viva saeta agarró él el primero.\\
Plantado ya en el umbral, al arco echábale tiento:\\
no lo tensó, que agotó antes sus manos en el tiramiento\\
---sus manos suaves, no usadas---, y dijo a sus compañeros:\\
«Amigos, no puedo armarlo, así que lo tome ya otro».
\end{versocita}

Como Leodes tiene dotes de adivino, más les habría valido a los pretendientes hacerle caso cuando dice:

\begin{versocita}
«A muchos de los mejores este arco les va a dejar rotos\\
el corazón y el aliento, pues es preferible del todo\\
morir a truncos vivir por no ver consumado ese logro\\
que aquí esperando y en corro nos tuvo un día tras otro.\\
Habrá quien guarde en su pecho aún esperanza, afanoso\\
por de Penélope ser, la mujer de Odiseo, el esposo.\\
Mas cuando al arco se atreva y ya lo tenga a sus ojos,\\
mejor que a alguna otra aquea veligalana obsequioso\\
por cortejarla se vuelva; y deje que al matrimonio\\
llame la reina al que, a más de fatal, venga más dadivoso».
\end{versocita}

Antínoo, después de burlarse de Leodes, ordena a Melantio que prepare un fuego para untar el arco con grasa y así ablandarlo. El cabrero obedece, pero ni aun así consigue ejecutar la proeza ninguno de la alegre banda. Solo quedan por probar Antínoo y Eurímaco. Entre tanto, Odiseo se escapa al patio y alista a sus dos leales, el boyero Filetio y el porquerizo Eumeo, y les da instrucciones para consumar su plan:

\begin{versocita}
«Ahora hay que entrar, y se hará uno por uno, no todos a un tiempo;\\
Yo iré primero, vosotros después; y a esta seña id atentos:\\
ninguno de esos galanos de los del ilustre cortejo\\
permitirá que la aljaba y el arco me den para el tiento;\\
pero, por toda la sala el arco llevándolo, Eumeo,\\
me lo pondrás en las manos, diciendo a las siervas tú luego\\
que echen cerrojo a las bien firmes puertas del gran aposento,\\
y si una escucha gemido o estruendo de hombres adentro\\
de nuestros muros, no vaya a querer asomarse allí abriendo,\\
no, que a su jera y labores habrá ella de estar en silencio.\\
Y a ti las puertas del patio te encargo, divino Filetio:\\
ciérralas con el cerrojo echándoles nudo de recio».
\end{versocita}

En la sala, Eurímaco fracasa con el arco y se queja amargamente.

\begin{versocita}
«Áyayme, a fe es por mí mismo y por todos la pena que siento,\\
y, aunque me amarga, no es tanto la boda lo que ahora lamento\\
---hay otras muchas aqueas y sin salir de este suelo en\\
Ítaca milsalinera, y muchas habrá en otros pueblos---,\\
como el que estemos en fuerza quedando todos tan lejos\\
del pardediós Odiseo; ni dar a la cuerda podemos\\
su arco: afrenta que oirán mañana los venideros».
\end{versocita}

El detalle no tiene desperdicio. A Eurímaco le preocupa más perder su reputación y que se sepa que ni él ni nadie de la pandilla han podido tensar el arco, que la boda con Penélope. Antínoo, por su parte, ni siquiera lo intenta y echa balones fuera comparando a Odiseo con un dios y, por tanto, dando a su fuerza carácter sobrenatural. Propone, a cambio, seguir con la fiesta. Cuando todos están bastante borrachos ---una circunstancia importante, ya que es más fácil abatir al enemigo si está ebrio, sin olvidar que el héroe ya ha aplicado ese truco con el cíclope---, Odiseo pone en marcha su treta:

\begin{versocita}
«Oíd, los que pretendéis con la bien clara reina alianza,\\
que esto es lo que el corazón desde sus entretelas me manda:\\
a Eurímaco sobre todo y a Antínoo divinaestampa\\
suplico ---y más cuando ha dicho a punto y verdad su palabra---\\
que ahora descansen del arco, que ya los dioses se encargan,\\
y al alba el dios le dará la fuerza a quien se le plazca.\\
Pero dejadme a mí el arco pulido y que prueba yo haga\\
ante vosotros de manos y fuerzas por ver si aún arraiga\\
aquel vigor que en mis miembros retuertos antaño llevaba,\\
o si me lo han ya matado el abandono y la errancia».
\end{versocita}

A los pretendientes la propuesta les sienta mal, entre otras cosas por miedo a que un mendigo sea capaz de triunfar donde ellos han fracasado, y Antínoo le increpa:

\begin{versocita}
«Ah, condenado extranjero, en ti no hay ni miaja de seso;\\
¿no te contentas con mesa aquí haber con nosotros, ileso,\\
entre tan alta ralea, sin mengua de yanta, y oyendo\\
nuestra conversa y palabra además? Ningún otro extranjero\\
palabra nuestra llegó aquí a escuchar, ni ningún pordiosero».
\end{versocita}

Luego le echa una diatriba, lo llama borracho y le amenaza:

\begin{versocita}
«Desgracia tanta en ti ahora, si el arco tensas, descuento,\\
pues que de nadie tendrás benevolencia y aprecio\\
en nuestro pueblo, que presto sin más en un negro velero\\
al rey Equeto, de los morideros verdugo y flagelo,\\
te enviaremos: de allí no hay regreso. Así que sereno\\
bebe, y porfía no busques con los que te son más mozuelos».
\end{versocita}

A lo que contesta Penélope:

\begin{versocita}
«No es noble, Antínoo, ni justo dejar de lo suyo privado\\
a quien llegó a este casal y ganó de Telémaco amparo.\\
¿Crees que si el huésped lograra tensar de Odiseo el gran arco\\
porque en su fuerza fiara y en el vigor de sus manos\\
me iba a llevar a su casa y a hacerme su esposa por pago?\\
A fe que ni él en su pecho está tal cosa esperando.\\
Ninguno se ande por eso hondodolido en su ánimo\\
mientras aquí banquetea, que no, que tampoco es el caso».
\end{versocita}

Eurímaco mete baza para explicar que ninguno de ellos ---todos hijos de nobles casas--- puede contemplar la posibilidad de que un mendigo arruine su reputación. Penélope le responde, sin pelos en la lengua:

\begin{versocita}
«¿Qué buena fama va a haber, Eurímaco, entre la gente\\
para quien casa deshonra comiendo en desprecio los beches\\
de su varón principal? Y por sobra os tomáis que se cuente.\\
El forastero es muy alto, y bien formado y es fuerte,\\
y de hijo ser de buen padre se precia por su progenie.\\
Así que dadle ya el arco pulido por ver si lo tiende.\\
Y cosa más aún os digo que cumpliré si sucede:\\
si lo templara, porque eso a él Apolo por gloria le diere,\\
lo vestiré, con jubón y manteo, de muy buenas vestes,\\
y le daré aguda azcona que a perros y hombres alueñe,\\
y espada doblefileña, y sandalias de buenos cordeles,\\
y su partida armaré adonde el corazón a ir lo lleve».
\end{versocita}

No pierdas detalle, biensabida lectora, bienatento lector. A diferencia de lo que nos cuenta Nolan, \emph{Penélope no sabe que el mendigo es su marido}. Y esto es importante porque su intervención favorece los planes de Ulises sin ella saberlo y porque muestra su generosidad y, sobre todo, \emph{su agencia}. La Penélope de Nolan ayuda a su esposo, al que ya ha reconocido, es decir, se comporta como su comparsa. La de Homero tiene voz y voluntad propias, piensa y decide por su cuenta. ¿Es nuestro bardo más feminista que el bienintencionado director?

A Telémaco, en cambio, le falta tiempo para sacar sus galones de hombre de casa:

\begin{versocita}
«Ninguno de los aqueos mejor que yo, madre mía,\\
para ofrecer o negar ese arco a quienquier que yo diga,\\
ni de los que señorean por los pedrizales de Ítaca,\\
ni de las islas que miran a Élide belyeguariza;\\
fuerza ninguno me hará malmigrado ni si es gana mía\\
dar a mi huésped el arco y que se lo lleve en propina.\\
Así que ve a tu aposento y de tu labor allí cuida,\\
que es el telar y la rueca, y di a tus mujeres que sigan\\
a su trabajo aplicadas, que el arco cuenta es masculina,\\
muy cuenta mía además, que soy yo quien la casa administra».
\end{versocita}

Muy astuto y ya es la segunda vez que nos lo muestra. De un solo golpe consigue mandar a su madre fuera de la estancia ---alejándola así de la inminente carnicería--- y hacerse con el arco, ya sabemos con qué propósito. La pobre Penélope, que no se espera que el hijo querido le meta semejante riña, se marcha apenada y echando de menos al marido ausente, hasta que Atenea, siempre atenta, la deja dormida:

\begin{versocita}
Ella marchó a su aposento otra vez, hondamente asombrada,\\
pues las palabras del hijo se entraron sagaces al alma.\\
Ya en las estancias de arriba, adonde subió con sus amas,\\
llora a Odiseo, a su esposo, y lo llora hasta que a ella le lanza\\
bajo los párpados, dulce, el sueño Atenea ojiglauca.
\end{versocita}

Eumeo, que, como sabemos, está en el ajo, se hace con el arma mientras los pretendientes se burlan de él. Luego Telémaco le pide a la nodriza Euriclea que eche el cerrojo a las puertas de la sala.

Entre tanto, Odiseo tensa el arco:

\begin{versocita}
Decían los pretendientes; en tanto el ardido Odiseo,\\
bien sopesado y mirado el gran arco por todo su cuerpo,\\
igual que el que hace del canto y de la cítara empleo\\
tensa en la nueva clavija la cuerda muy sin esfuerzo\\
luego de atar una tripa torcida de oveja a su extremo,\\
así, sin pena ni afán, tensó el gran arco Odiseo.\\
Tomó después, por probarlo, entre sus dedos el nervio,\\
y agudo y claro a su mano cantó, como voz de vencejo.
\end{versocita}

¡Qué maravilla! Ulises hace cantar el arco «como voz de vencejo» y a sir Christopher ese detalle no se le escapa y, de hecho, lo reviste de significado: su héroe hace sonar el arco \emph{siempre} para advertir a sus víctimas de lo que les espera y darles una oportunidad. El del Poeta, por el contrario, carece por completo de escrúpulos.

Antes de la matanza, el ritual que lo identifica a los ojos de todos:

\begin{versocita}
Tomó una vívida flecha que sobre la mesa y en huelgo\\
nuda yacía; guardaba la aljaba en su cuenco agujero\\
a las demás, las que pronto probaran allí los aqueos.\\
Y ya, la vira en su muesca, apuntando tiraba del nervio\\
desde su escaño sentado, y el dardo lanzó, que a derecho\\
pasó por todas las hachas sin yerro alguno en su vuelo\\
de ojo en ojo a través, hasta haber en el patio su término\\
el tiro broncicebado. Y entonces dijo a Telémaco:
\end{versocita}

\begin{versocita}
«Telémaco, no te avergüenza el huésped que tienes sentado\\
aquí a tu mesa; ni el blanco lo erré ni estirar ese arco\\
largo cansancio me dio: está firme aún el brío en mi brazo,\\
no como, a mi menosprecio, tasábanmelo estos galanos.\\
Hora es de darles ahora también la cena a los dánaos\\
mientras hay luz, que vendrá otro solaz enseguida a animarnos:\\
baile y canción, y el festín en sazón quedará y adornado».
\end{versocita}

Humor negro no le falta. Todos sabemos qué cena les espera a los dánaos:

\begin{versocita}
Dijo, y de cejas dio seña; y está aguda espada ciñendo\\
ya el de Odiseo del cielo, el hijo querido, Telémaco,\\
y, echando mano a su lanza, junto al escañuelo paterno\\
se planta muy bien armado de un bronce como de fuego.
\end{versocita}

Nada de «tú quieto, hijito, que papá se ocupa». Telémaco es un héroe por derecho propio: quiere defender a su familia, vengar las afrentas recibidas y recuperar su hacienda, igual que su padre. No necesita este guerrero, «armado de un bronce como de fuego», que nadie le ofrezca un \emph{safe space}.
